% Einheit 5 von 5 – Körper im Koordinatensystem und Drehungen (bank/punkte-und-strecken-im-koordinatensystem/e5.jsonl, zusammenbau v0.1)
%% TODO Zeitmarke Einheit 5: Marken-Zeile nicht lesbar
%% TODO Zweigzeile Teil 1: was hier gelernt wird (Satz in Schülersprache aus dem Einheitstitel) – Einheit 5
\einheitenkopf[e5]{Einheit 5 von 5 · Körper im Koordinatensystem und Drehungen}
\zweigzeile{Abitur GK}
\setcounter{aufgabe}{18}

%% TODO Titel: Ich-kann-Satz fehlt in der Bank (Platzhalter: Kettenname) – Nr. 19
%% TODO Anweisung: ein Satz über der ersten Teilaufgabe fehlt in der Bank – Nr. 19
\begin{aufgabe}{Körper und Drehungen}
%% TODO \rechenplatz{n}: Schrittzahl des Grundfalls fehlt in der Bank (nur bei fester Schreibform, 2.3 b)
\begin{teile}
\teil Ein Quader steht auf der $x_1x_2$-Ebene; seine Deckenecke $G$ hat $x_3 = 4$. Was sagt das über die Deckfläche? \\ \kreuz{Sie liegt in der $x_1x_2$-Ebene.} \\ \kreuz{Sie liegt parallel zur $x_1x_2$-Ebene in der Höhe $4$.} \\ \kreuz{Sie steht senkrecht auf der $x_1x_2$-Ebene.}
\teil Gerades Prisma $ABCDEF$ ($D$ über $A$, $E$ über $B$, $F$ über $C$) mit $A(2 | 1 | 0)$, $B(7 | 1 | 0)$, $C(2 | 4 | 0)$, $D(2 | 1 | 3)$: $F$? \hfill (Abitur 2026 GK) \\ F( \leerfeld | \leerfeld | \leerfeld )
\teil Eine Pyramide hat als Grundfläche ein rechtwinkliges Dreieck mit den Katheten $6$ cm und $8$ cm; die Spitze liegt $5$ cm senkrecht über der Ecke mit dem rechten Winkel. Wähle geeignete Koordinaten der vier Ecken (1 LE = 1 cm). \hfill (Abitur 2020 GK)
\teil Ein Würfel mit Kante $2$ hat eine Ecke im Ursprung, alle Ecken haben Koordinaten $\ge 0$. Er wird um $(-1 | -1 | 1)$ verschoben. Welche Ecke hat danach $x_1 < 0$, $x_2 > 0$ und $x_3 > 2$? \hfill (Abitur 2024 GK) \\ ( \leerfeld | \leerfeld | \leerfeld )
\teil Eine Pyramide hat die Grundfläche $A(4 | -1 | 2)$, $B(1 | 3 | 2)$, $C(-1 | 0 | 2)$ und die Spitze $S(2 | 1 | 7)$. Begründe: Die Grundfläche ist parallel zur $x_1x_2$-Ebene. Höhe der Pyramide? \hfill (Abitur 2018 GK) \\ h = \leerfeld
\teil Ein gerader Zylinder mit Radius $5$ und Höhe $6$ steht auf der $x_1x_2$-Ebene; $M(4 | 3 | 6)$ ist der Mittelpunkt der Deckfläche. Prüfe, ob $P(7 | 7 | 0)$ auf dem Rand der Grundfläche liegt. \hfill (Abitur 2020 LK)
\teil Ein Dachgeschoss hat einen rechteckigen Boden von $10$ m mal $8$ m. Die Dachflächen steigen von den beiden Traufwänden (Höhe $0{,}5$ m) gleichmäßig zum First über der Mitte (Höhe $5{,}5$ m). Flächen mit weniger als $2$ m Raumhöhe zählen nicht zur Wohnfläche. Welcher Anteil des Bodens zählt nicht? \hfill (Abitur 2019 GK) \\ \leerfeld \%
\teil Der Punkt $B(0 | 10 | 1)$ wird um die $x_3$-Achse gedreht. Gib drei Bildpunkte an, bei denen eine Koordinate $6$ ist. \hfill (Abitur 2024 LK)
\teil Das Dreieck $A(5 | 3 | 0)$, $B(3 | 5 | 0)$, $C(3 | 3 | 4)$ wird um die Kante $AB$ gedreht, bis $C$ in der $x_1x_2$-Ebene liegt, mit $x_1 < 4$. Koordinaten des gedrehten Punktes $C'$? \hfill (Abitur 2023 GK) \\ C'( \leerfeld | \leerfeld | \leerfeld )
\end{teile}
\end{aufgabe}

%% TODO Titel: Ich-kann-Satz fehlt in der Bank (Platzhalter: Kettenname) – Nr. 20
%% TODO Anweisung: ein Satz über der ersten Teilaufgabe fehlt in der Bank – Nr. 20
\begin{aufgabe}{Körper und Drehungen – Fehler finden, Begründen}
\begin{teile}
\teil Ich finde den Fehler: Gerades Prisma $ABCDEF$ ($D$ über $A$, $E$ über $B$, $F$ über $C$) mit $A(1 | 2 | 0)$, $B(6 | 2 | 0)$, $C(1 | 5 | 0)$ und der Höhe $4$. Paul: $F(6 | 2 | 4)$.
\teil Begründe: Beim geraden Prisma auf der $x_1x_2$-Ebene hat jede Deckenecke dieselbe erste und zweite Koordinate wie ihre Grundecke.
\end{teile}
\end{aufgabe}
